\subsection*{Цель работы}
Ознакомиться с основными приемами работы с документацией при составлении программ для микроконтроллеров.

\subsection*{Постановка задачи}
Создать новый проект в Keil $\mu$Vision5 и разработать программу для микроконтроллера (МК) STM32F200, которая включает и выключает светодиоды.

\subsection*{Методика выполнения работы}

Изначально была реализована простейшая программа, управляющая одним мигающим светодиодом. Это позволило разобраться с принципом работы с регистрами GPIO.

\subsection*{Программа 1: мигание одним светодиодом}

Первым шагом необходимо включить тактирование порта GPIOG, записав бит 6 в регистр \texttt{RCC\_AHB1ENR}:

\begin{lstlisting}[language=C]
*(unsigned long*)(0x40023830) |= 0x40;
for(i=0; i<4; i++){}
\end{lstlisting}

Пустой цикл после записи обеспечивает короткую задержку, необходимую для стабилизации работы периферийного блока.

Далее настраивается режим вывода PG7. В регистре \texttt{GPIOG\_MODER} пара битов [15:14] управляет режимом пина PG7. Значение \texttt{01} соответствует режиму вывода.

\begin{lstlisting}[language=C]
*(unsigned long*)(0x40021800) = (*(unsigned long*)(0x40021800) &
(~0x00008000)) | (0x00004000);
\end{lstlisting}

Здесь сначала очищается бит 15 (маска \texttt{\textasciitilde 0x00008000}), затем устанавливается бит 14 (маска \texttt{0x00004000}).

В основном цикле светодиод попеременно включается и выключается через регистр \texttt{GPIOG\_ODR} с программной задержкой между переключениями:

\begin{lstlisting}[language=C]
*(unsigned long*)(0x40021814) |= 0x80;    // PG7 on
for( j=0; j<2000000 ;j++ ){}
*(unsigned long*)(0x40021814) &= ~0x80;   // PG7 off
for(j=0; j<2000000 ; j++){}
\end{lstlisting}

Установка бита 7 в \texttt{GPIOG\_ODR} (адрес \texttt{0x40021814}) включает светодиод, сброс~--- выключает. Задержка в 2\,000\,000 итераций обеспечивает видимое мигание.

На фотографиях \ref{fig:2026-02-25_16.07.24.jpg} и \ref{fig:2026-02-25_16.08.37.jpg} приведён пример мигания светодиода.

\screenshot{2026-02-25_16.07.24.jpg}{Светодиод PH7 выключен.}
\screenshot{2026-02-25_16.08.37.jpg}{Светодиод PH7 включен.}

\subsection*{Программа 2: управление несколькими светодиодами}

После освоения работы с одним светодиодом была реализована программа, управляющая восемью светодиодами на портах G, H и I, создающая эффект бегущего огня.

\textbf{Включение тактирования портов G, H, I.}

\begin{lstlisting}[language=C]
*(unsigned long*)(0x40023830) |= (0x40 | 0x80 | 0x100);
\end{lstlisting}

Биты 6, 7 и 8 регистра \texttt{RCC\_AHB1ENR} включают тактирование GPIOG, GPIOH и GPIOI соответственно.

\textbf{Настройка выводов PG6, PG7, PG8 как выходов.}

\begin{lstlisting}[language=C]
*(unsigned long*)(0x40021800) &= ~((3 << 12) | (3 << 14) | (3 << 16));
*(unsigned long*)(0x40021800) |=  ((1 << 12) | (1 << 14) | (1 << 16));
\end{lstlisting}

В регистре \texttt{GPIOG\_MODER} очищаются пары битов [13:12], [15:14], [17:16], а затем в каждой паре устанавливается значение \texttt{01}~--- режим вывода.

\textbf{Настройка выводов PH2, PH3, PH6, PH7 как выходов.}

\begin{lstlisting}[language=C]
*(unsigned long*)(0x40021C00) &= ~((3 << 4) | (3 << 6)
                                 | (3 << 12) | (3 << 14));
*(unsigned long*)(0x40021C00) |=  ((1 << 4) | (1 << 6)
                                 | (1 << 12) | (1 << 14));
\end{lstlisting}

Аналогичная операция для регистра \texttt{GPIOH\_MODER}: очистка и установка пар битов для выводов PH2, PH3, PH6 и PH7.


\textbf{Основной цикл~--- четыре фазы бегущего огня.}

Всё, что ниже, находится в постоянно исполняющемся цикле \verb|while(1) {... |, соответственно, в конце следует закрывающая скобка \verb|...}|.

\textbf{Фаза 1:} включаются PH3 и PH6.

\begin{lstlisting}[language=C]
*(unsigned long*)(0x40021C14) |= (0x08 | 0x40);   // PH3, PH6 on
for (j = 0; j < 4000000; j++) {}
*(unsigned long*)(0x40021C14) &= ~(0x08 | 0x40);  // PH3, PH6 off
\end{lstlisting}

В регистре \texttt{GPIOH\_ODR} (адрес \texttt{0x40021C14}) устанавливаются биты 3 и 6 (маски \texttt{0x08} и \texttt{0x40}), что включает светодиоды на выводах PH3 и PH6. После программной задержки в 4\,000\,000 итераций эти же биты сбрасываются операцией AND с инвертированной маской, и светодиоды выключаются.


\textbf{Фаза 2:} включаются PH7 и PI10.

\begin{lstlisting}[language=C]
*(unsigned long*)(0x40021C14) |= 0x80;   // PH7 on
*(unsigned long*)(0x40022014) |= 0x400;  // PI10 on
for (j = 0; j < 4000000; j++) {}
*(unsigned long*)(0x40021C14) &= ~0x80;  // PH7 off
*(unsigned long*)(0x40022014) &= ~0x400; // PI10 off
\end{lstlisting}

Здесь задействованы два разных порта одновременно: бит 7 (маска \texttt{0x80}) устанавливается в \texttt{GPIOH\_ODR}, а бит 10 (маска \texttt{0x400})~--- в \texttt{GPIOI\_ODR} (адрес \texttt{0x40022014}). Это включает светодиоды PH7 и PI10. После задержки оба бита сбрасываются в соответствующих регистрах.

\textbf{Фаза 3:} включаются PG6 и PG7.

\begin{lstlisting}[language=C]
*(unsigned long*)(0x40021814) |= (0x40 | 0x80);   // PG6, PG7 on
for (j = 0; j < 4000000; j++) {}
*(unsigned long*)(0x40021814) &= ~(0x40 | 0x80);  // PG6, PG7 off
\end{lstlisting}

В регистре \texttt{GPIOG\_ODR} (адрес \texttt{0x40021814}) устанавливаются биты 6 и 7 (маски \texttt{0x40} и \texttt{0x80}), включая светодиоды на выводах PG6 и PG7. После задержки биты сбрасываются.

\textbf{Фаза 4:} включаются PG8 и PH2.

\begin{lstlisting}[language=C]
*(unsigned long*)(0x40021814) |= 0x100;  // PG8 on
*(unsigned long*)(0x40021C14) |= 0x04;   // PH2 on
for (j = 0; j < 4000000; j++) {}
*(unsigned long*)(0x40021814) &= ~0x100; // PG8 off
*(unsigned long*)(0x40021C14) &= ~0x04;  // PH2 off
\end{lstlisting}

Снова задействованы два порта: бит 8 (маска \texttt{0x100}) в \texttt{GPIOG\_ODR} включает PG8, а бит 2 (маска \texttt{0x04}) в \texttt{GPIOH\_ODR} включает PH2. После задержки оба светодиода выключаются, и цикл начинается заново с первой фазы.

Между фазами выдерживается программная задержка длительностью 4\,000\,000 итераций, что обеспечивает визуально различимое переключение светодиодов. После четвёртой фазы цикл повторяется.

На фотографиях \ref{fig:2026-02-25_16.08.42.jpg}~-- \ref{fig:2026-02-25_16.08.58.jpg} показана работа программы на отладочной плате. Видно, как различные группы светодиодов зажигаются поочерёдно в соответствии с четырьмя фазами основного цикла.

\screenshot{2026-02-25_16.08.42.jpg}{Первая фаза: светодиоды PH3 и PH6 включены}
\screenshot{2026-02-25_16.08.48.jpg}{Вторая фаза: светодиоды PH7 и PI10 включены}
\screenshot{2026-02-25_16.08.52.jpg}{Третья фаза: светодиоды PG6 и PG7 включены}
\screenshot{2026-02-25_16.08.58.jpg}{Четвёртая фаза: светодиоды PG8 и PH2 включены}

\subsection*{Результаты работы}

Таким образом, в ходе выполнения лабораторной работы был разработан проект в среде Keil $\mu$Vision5 и написана программа на языке C для микроконтроллера STM32F200.

Сначала была реализована простейшая программа, управляющая мигающим светодиодом на выводе PG7, что позволило освоить базовые операции: включение тактирования порта, настройку вывода как выхода и управление его состоянием. Затем программа была расширена до управления восемью светодиодами на трёх портах (G, H, I), реализующего эффект бегущего огня с четырьмя фазами переключения.
